\documentclass{article}

\usepackage[margin=1in]{geometry}
\usepackage{amsmath}
\usepackage{graphicx}
\usepackage{booktabs}
\usepackage{hyperref}
\usepackage{float}
\usepackage{tikz}
\usepackage{pgfplots}
\usepackage{pgfplotstable}
\usepackage{placeins}
\usepackage[most]{tcolorbox}
\usepackage[all]{hypcap}

\pgfplotsset{compat=1.18}


\title{Analysis of Parallelization Techniques on k-mer Counting for Genomic Sequences}

\author{Alex McGinn}
\date{\today}

\begin{document}

\maketitle

%========================================================== Introduction
\section{Introduction}

\vspace{\baselineskip}

In genomic sequence analysis, genetic information is represented by long sequences containing millions or even billions of nucleotide 
bases, primarily adenine (A), cytosine (C), guanine (G), and thymine (T). To compare genomic samples, scientists need methods that can 
identify meaningful similarities even when two sequences are not exact matches. This is important because even closely related organisms 
or samples can contain mutations, insertions, deletions, and other sequence variations. Although genomic sequences may differ, 
biologically meaningful similarities can often be detected through shared local sequence patterns such as k-mers \cite{moeckel2024kmer}.

\vspace{\baselineskip}

One method for identifying these local patterns is k-mer analysis. A sliding window of length \(k\) is moved across a genomic sequence, 
producing every contiguous substring of length \(k\). These substrings are called k-mers, and the number of times each unique k-mer occurs 
can then be counted. The resulting k-mer frequencies provide a compact way to characterize and compare genomic sequences without requiring 
the sequences to match in their entirety. Larger values of \(k\) provide more sequence-specific information, while smaller values provide 
greater sensitivity to local similarities \cite{lees2016seer}. For example, if two genomic regions contain the same 15-mer, this represents 
a more specific local match than sharing a much shorter k-mer. However, a single base difference can disrupt several overlapping large k-mers, 
while shorter k-mers are less sensitive to an individual mismatch.

\vspace{\baselineskip}

When analyzing a genetic sample, we do not obtain the entire genome in one clean pass. Instead, sequencing produces many shorter reads from the sample,
which often overlap and contain repeated coverage of the same genomic regions. As a result, the same k-mer may appear several times across the reads.
Because the complete genome may not already be known, we cannot always determine its total size directly from the sequencing output. However, as Marcais
and Kingsford explain, k-mer frequency statistics can be used to estimate sequencing coverage and subsequently ``derive an estimate of the genome size
from (c) and the total length of the reads'' \cite{marcais2011jellyfish}.

\vspace{\baselineskip}

Although, as k increases, the number of possible k-mers grows exponentially, leading to an increased amount of memory required to 
store and process k-mer counts. Alternatively, if k is shorter, we may have many matches, which would require the counter variable to 
have a larger range. 
For large genomic datasets, these computations become expensive, but there are various methods to accelerate this 
process using parallelization. The aim of this project is therefore to analyze the effectiveness of CPU/GPU parallelism for these  
types of tasks.

%========================================================== Background
\section{Background}

\vspace{\baselineskip}

To process a set of k-mer (determined by the k value; e.g. k = 4) we first move a sliding window of size k across the entire sequence and extract every k-mer substring.  
The first benefit is that we can now compare genomic sequences without simply comparing their complete ordering. Another benefit is that when we 
break a DNA sequence into many smaller subsequences of a known size, the condensed format allows us to perform computations in parallel, 
leading to a needed speedup. Although, as the input dataset size and k-value increase, the ability to store the data necessary in main memory becomes 
increasingly difficult.  

\vspace{\baselineskip}

To address this problem, one strategy for storing k-mers is to index an array or a similar data structure for each possible sequence. 
This allows us to predict the amount of main memory necessary to store all possible combinations for a dataset and set k-mer value.  
For example, a dataset could have a dedicated array that contains all possible k-mer permutation for k-value set equal to five.  
We know whether this would be possible or not because the number of possible permutations of such a string with 4  
possible bases is:
 
\[
4^k
\]

e.g. in a 5 k-mer example:

\[
4^5 = 1024 \text{ possible subsequences}
\]

\vspace{\baselineskip}

We could then also have an array that holds the count of each k-mer in the dataset. Therefore, we know that we  
need 1024 representations of the four possible bases. We can then also predict the size of the array used to count those  
occurrences. To determine the size necessary, we then need to consider the maximum possible number of occurences  
a k-mer can have. This will be dictated by frequency, which we cannot outright predict. However, we  
can be certain of the total number of possible occurrences based on the dataset size. That is because the maximum  
value will always be equal to how many possible substrings there are, determined by:

\vspace{\baselineskip}

\[
n - k + 1
\]

\vspace{\baselineskip}

e.g. in a DNA Sequence of length 66,000 with k-value equal to 5:

\vspace{\baselineskip}

\[
66000 - 5 + 1 = 65996
\]

\vspace{\baselineskip}

Which means that we need array elements that go up to at least 65996. In this instance, we would know we now need  
elements larger than an unsigned short, which would reach capacity at 65,535

\vspace{\baselineskip}

\subsection{Related Work}

\vspace{\baselineskip}

One similar project proposed by Marcais and Kingsford \cite{marcais2011jellyfish} introduced Jellyfish, a parallel k-mer counting 
algorithm designed to improve performance and memory efficiency when processing large genomic datasets. Jellyfish utilizes a
multithreaded, lock-free hash table that allows concurrent updates through atomic operations rather than conventional
locking mechanisms. The implementation also leverages compact data representations to reduce memory requirements. Their work demonstrates 
the importance of synchronization overhead and memory-efficient data structures in parallel k-mer counting, both of which were 
considered when designing and evaluating the parallel implementations in this project.

\vspace{\baselineskip}

Deorowicz et al. \cite{deorowicz2015kmc2} introduced KMC 2, a parallel k-mer counting algorithm that focuses on reducing memory 
consumption and disk I/O overhead when processing large genomic datasets. Unlike approaches that rely primarily on concurrent hash-table 
updates, KMC 2 partitions sequence data into smaller groups using signatures derived from minimizers. These partitions can then be processed 
using parallel sorting and counting operations, reducing the amount of data that must be held in memory simultaneously. The authors also 
introduced $(k,x)$-mers to reduce redundant processing of overlapping k-mers. This approach highlights the importance of memory management, 
workload distribution, and I/O overhead, which are relevant considerations for the parallelization methodology implemented in this project.

\vspace{\baselineskip}

Erbert et al. \cite{erbert2017gerbil} introduced Gerbil, a parallel k-mer counting tool designed
to improve performance and memory efficiency, particularly for larger k-mer lengths.
Gerbil uses a two-stage approach that first partitions genomic sequence data into temporary
files, then processes the resulting k-mers using hash tables. Unlike implementations that
allocate a frequency array containing all $4^k$ possible k-mers, Gerbil uses hash tables to
store the observed k-mers, allowing it to support larger k-mer lengths without requiring
an exponentially sized dense frequency array.

\vspace{\baselineskip}

Gerbil also supports CUDA acceleration by distributing k-mer counting operations across GPU
threads. The implementation maintains hash tables in GPU memory and uses atomic operations
to coordinate concurrent updates. The authors emphasize the importance of memory access
patterns, workload distribution, and GPU memory limitations when designing an efficient
parallel counting algorithm. These considerations are relevant to this project, where CUDA
uses atomic increments to update a shared frequency array. However, unlike Gerbil's
hash-table approach, the CUDA implementation developed in this project uses a direct-indexed
frequency array. This simplifies counter access but limits the practical k-mer length due
to the exponential growth in memory requirements.

%========================================================== Sequential Implementation

\section{Sequential Implementation}

\vspace{\baselineskip}

\subsection{Dataset and FASTA Input Processing}

\vspace{\baselineskip}

The initial sequential implementation was developed in C++ and tested using the Escherichia coli 
K-12 MG1655 reference genome, obtained from the National Center for Biotechnology Information (NCBI).
The reference genome contains 4,641,652 nucleotide bases and is provided in FASTA format \cite{ncbi_ecoli}.

\vspace{\baselineskip}

To evaluate performance on a larger genomic dataset, human chromosome 18 from the GRCh38
reference genome was also used \cite{ucsc_chr18}. The chromosome contains approximately 80.4 million nucleotide
positions, making it substantially larger than the E. coli reference genome. In addition,
the chromosome contains ambiguous nucleotide positions represented by N, providing an
opportunity to evaluate how the implementations handle invalid k-mer windows. Both datasets
were used to benchmark the sequential, OpenMP, and CUDA implementations.

\vspace{\baselineskip}

FASTA files contain nucleotide bases separated by header lines beginning with the `\texttt{>}' character.
The implementation reads the file line by line, concatenating nucleotide bases belonging to the same sequence
and storing each completed sequence in a string.

\vspace{\baselineskip}

The runtime is evaluated both before the FASTA loading operation is performed and before the sequential counting benchmark begins, allowing the
execution time of the counting algorithm to be measured independently of file input operations. 

\vspace{\baselineskip}

\subsection{Base Encoding}

\vspace{\baselineskip}

Each base is represented using a two-bit integer encoding, where A, C, G, and T correspond to the values
0, 1, 2, and 3, respectively. Since each base requires two bits, a k-mer of length $k$ can be represented
using $2k$ bits. The implementation then stores the encoded k-mer in an unsigned 64-bit integer using bit manipulation, 
which allows a maximum k-mer length of $k=32$. However, the dense frequency array requires $4^k$ counters, making memory 
consumption the primary limitation on the supported k-mer length rather than the integer encoding itself.

\vspace{\baselineskip}

A count \texttt{vector} is allocated with $4^k$ elements, corresponding to every possible k-mer of length $k$.
Each element is a \texttt{uint64\_t}, resulting in a total allocation of $8 \times 4^k$ bytes. As each
k-mer is processed, its two-bit representation is shifted into the encoded integer. The resulting integer
value is then used directly as the index into the count \texttt{vector}, allowing the corresponding k-mer count
to be incremented without requiring a hash-table lookup.

\vspace{\baselineskip}

Many genetic sequences will have bases that are not read correctly, which is represented in the file as an \texttt{N}. 
Entries that contain this character representing an invalid reading will not be used in the final count.

\subsection{Sliding Window and Frequency Counting}

The sequential algorithm processes each sequence using a sliding window of length $k$. Rather than
constructing a new substring at every position, the implementation maintains an integer containing the
encoded bases of the current window.

\vspace{\baselineskip}

For each new base, the encoded integer is shifted left by two bits, and the encoded value of the new
nucleotide is inserted using a bitwise OR operation. A bitmask is then applied to remove any bits extending
beyond the current k-mer, ensuring that only the most recent $k$ bases are retained.

\vspace{\baselineskip}

Once $k$ valid nucleotides have been processed, the resulting integer is used as an index into a frequency
counting array. Each array entry stores the number of occurrences of its corresponding k-mer. This allows
the frequency counts to be updated directly without requiring a hash table or additional lookup operations.

\vspace{\baselineskip}

The accuracy of the output was verified against an independent Python implementation that
counts k-mer frequencies using string operations rather than the two-bit encoding method.

\vspace{\baselineskip}

\subsection{Time \& Memory Complexity}

\vspace{\baselineskip}

The main components of the applications are reading the sequences from the file, initializing the counter array, and processing and incrementing each k-mer. 
Initializing the count of the array to default values takes $O(4^k)$ time, while processing and incrementing each k-mer only takes $O(n)$ time. Since reading 
from a file also takes linear time with respect to the number of bases, the overall time complexity is $O(4^k + n)$. This sequential implementation was 
evaluated for 3, 6, 9, and 12.

\vspace{\baselineskip}

The memory constraints for this implementation require us to hold the sequence itself in main memory, as well as a \texttt{vector} of \texttt{uint\_64t} 
that has enough memory for all $4^k$ possible sequences. Table \ref{tab:sequential_results} shows that at $k=12$ the amount of RAM that is necessary to store 
all possible sequences increases dramatically.

\vspace{\baselineskip}

\FloatBarrier
\input{/home/amcgin/kmerCounting/source/sequential/etc/benchmark_table.tex}
\FloatBarrier

\vspace{\baselineskip}

In Figure \ref{fig:sequential_performance}, we can see that the time it takes to process the same amount of data for different k-values also increases greatly 
after $k=9$. This is likely due to the size of the count array needed for all possible values. Initializing that array to default values alone is on the 
order of $4^{12}= 16{,}777{,}216$. With that many possible k-mers, the larger frequency array may reduce cache efficiency and increase memory access latency, 
which is likely why we see such a drop off in the performance, despite the sequence size remaining constant.

\vspace{\baselineskip}

\FloatBarrier
\input{/home/amcgin/kmerCounting/source/sequential/etc/benchmark_plot.tex}
\FloatBarrier

\vspace{\baselineskip}

\subsection{Correctness Verification}

\vspace{\baselineskip}

The sequential, OpenMP, and CUDA implementations were verified against an independent Python reference.
The implementations were tested using the E. coli reference genome and human chromosome 18 datasets.
For chromosome 18 at \(k=12\), both parallel implementations counted 80,089,048 valid k-mers, representing 
11,177,170 unique k-mers. The frequency of every unique k-mer matched the Python reference, producing zero 
mismatches for both implementations.

\vspace{\baselineskip}

Verification was performed for k-mer lengths of 3, 6, 9, and 12. For each configuration,
the frequency of every unique k-mer was compared against the Python reference. All 
configurations produced identical frequency distributions with zero mismatches, confirming
the correctness of each implementation for the tested configurations.

%========================================================== Parallel Methodology
\section{Parallelization Methodology}

\vspace{\baselineskip}

The objective of this project is to analyze the effectiveness of CPU/GPU parallelization on k-mer counting for genomic sequences. 
The algorithms were written using OpenMP for parallelization on multiple CPU cores and CUDA for GPU parallelization. 

\vspace{\baselineskip}


Both implementations were evaluated using the Escherichia coli reference genome and human chromosome 18, with k-mer lengths of 3, 6, 9, and 12. 
Performance was measured across ten trials per configuration. The evaluation considers counting time, speedup relative to the sequential baseline, and memory usage. 
OpenMP experiments examine thread scaling, thread-local frequency arrays, and merging overhead. CUDA experiments distinguish kernel execution time, warm-context counting time, 
and cold-start counting time, allowing GPU initialization overhead to be evaluated separately.

\vspace{\baselineskip}


%========================================================== OpenMP

\subsection{OpenMP CPU Parallelization}

\vspace{\baselineskip}

The initial implementation uses OpenMP to distribute the k-mer counting workload across multiple CPU threads. The input sequences are partitioned 
into smaller regions and then distributed across the threads to allow processing of different portions of the data. However, k-mers are generated using 
overlapping windows, so the implementation accounts for the necessary overlap in between partitions. 

\vspace{\baselineskip}

There are three approaches that can be taken when working with OpenMP and our current design. First, we could use the same count array and 
share it between the threads, but this requires some sort of protection to ensure thread safety. Atomics could be a good solution, and requires 
minimal changes to the sequential logic. Although, if many of the threads are attempting to increment the same counter, contention 
could negatively affect the performance impact. 

\vspace{\baselineskip}

Another method would be to assign each thread with its own array, and then consolidate the data after each thread finishes counting. This would 
likely lead to an increased problem with the amount of memory required for larger k, which was encountered in the sequential approach. This 
is because we would have to initialize a $4^k$ array for every thread. For 8 threads and $k=12$ this would already take up 1GB of space. 

\vspace{\baselineskip}

The last method would be partitioned ownership, in which each thread is designated owner of some part of the shared array. The data is distributed, but 
when a thread encounters a k-mer that does not belong to its partition, it instead determines the owner and sends the k-mer to their buffer. 

\vspace{\baselineskip}

Of the three approaches considered, the thread-local counter method was selected for the initial OpenMP implementation.
This approach avoids synchronization during frequency updates by providing each thread with an independent counter array.
Although this introduces additional memory overhead, it avoids contention on shared counters. The partitioned ownership approach remains a 
potential direction for future optimization, particularly for larger k-mer lengths where the memory requirements of thread-local arrays become significant.

\vspace{\baselineskip}

The implemented algorithm divides each input sequence into blocks containing up to 65,536 k-mer starting positions.
These blocks are distributed among CPU threads using OpenMP with static scheduling. Because adjacent k-mers overlap,
each thread may read up to k-1 additional nucleotide bases beyond its assigned block boundary, ensuring that all
valid k-mers are counted without duplicating starting positions.

\vspace{\baselineskip}

Each OpenMP thread maintains a private frequency array containing $4^k$ counters, initialized in parallel before
the counting phase begins. Threads process their assigned blocks independently using the same two-bit encoding
and sliding-window technique as the sequential implementation. This avoids race conditions because threads do
not update a shared frequency array during counting.

\vspace{\baselineskip}

Once counting is complete, the private frequency arrays are merged into the global frequency array. The merging
operation is parallelized by distributing ranges of counter indices among threads, ensuring that each output
counter is updated by only one thread. This eliminates the need for atomic operations during the merging phase.

\vspace{\baselineskip}

The primary limitation of this implementation is memory consumption. Since each thread maintains a separate
counter array, memory requirements increase approximately as $O(p4^k)$, where $p$ is the number of threads.
At $k=12$, each private array requires 128 MiB, resulting in approximately 1 GiB of private counter storage
when using eight threads.

\vspace{\baselineskip}

%========================================================== CUDA 

\subsection{CUDA GPU Parallelization}

\vspace{\baselineskip}

A second parallel implementation uses CUDA to perform k-mer counting on the GPU. Sequence data is transferred from RAM into the device memory 
of the GPU. The initial implementation assigns individual k-mer starting positions to GPU threads. Each thread encodes its assigned k-mers using the 
same two bit representation as the sequential implementation. 

\vspace{\baselineskip}

The counting kernel uses a one-dimensional grid of thread blocks, with each thread processing one
k-mer starting position. Each thread reads the following k nucleotide bases and constructs the
encoded value using bitwise operations. Windows containing invalid nucleotides are skipped. Unlike
the sequential implementation, which maintains a rolling encoding across consecutive positions,
CUDA threads independently encode their assigned k-mers. This allows the counting workload to be
distributed across many GPU threads.

\vspace{\baselineskip}

The encoded k-mers resulting from this task are used to index a frequency array stored in GPU global memory. This implementation deals with race conditions 
using atomic increments. 

\vspace{\baselineskip}

The CUDA implementation was evaluated by measuring the execution time of the counting kernel separately
from the overall counting operation. GPU initialization was also measured independently to distinguish
cold-start performance from warm-context performance. The warm-context measurement includes device
memory allocation, host-to-device and device-to-host transfers, kernel execution, and cleanup, while
the cold-start measurement additionally includes CUDA context initialization. Experiments evaluated
different k-mer lengths and dataset sizes using a fixed configuration of 256 threads per block.

\vspace{\baselineskip}

\subsection{CPU \& Compiler Info}

The sequential and OpenMP implementations were compiled using the GNU C++ compiler (g++) with the C++17
language standard and the -O2 optimization flag. This optimization level enables compiler
optimizations intended to improve execution performance while retaining the same algorithmic
behavior. OpenMP was also compiled using the -fopenmp compiler flag. The CUDA implementation was compiled using 
NVIDIA's CUDA compiler (nvcc) with C++17, the -O2 optimization flag, and the sm\_120 target architecture. All 
implementations were executed under AlmaLinux 9.6 within WSL2 on a Windows host. Each benchmark configuration 
was executed ten times to measure average execution time and the variation between trials.
\begin{itemize}
    \item \textbf{CPU:} AMD Ryzen 7 7700 8-Core Processor
    \item \textbf{GCC:} g++ (GCC) 11.5.0 20240719 (Red Hat 11.5.0-14)
    \item \textbf{Architecture:} x86\_64
    \item \textbf{Physical Cores:} 8
    \item \textbf{Hardware Threads:} 16 (2 per core)
    \item \textbf{L1 Data Cache:} 256 KiB total, 32 KiB per core
    \item \textbf{L1 Instruction Cache:} 256 KiB total, 32 KiB per core
    \item \textbf{L2 Cache:} 8 MiB total, 1 MiB per core
    \item \textbf{L3 Cache:} 32 MiB
    \item \textbf{Execution Environment:} WSL2 / Microsoft Hyper-V
\end{itemize}

\vspace{\baselineskip}

\FloatBarrier

\subsection{GPU Info}

\vspace{\baselineskip}

\begin{table}[H]
    \centering
    \caption{GPU hardware specifications for CUDA experiments}
    \label{tab:gpu_specs}
    \begin{tabular}{lc}
        \toprule
        \textbf{Specification} & \textbf{Value} \\
        \midrule
        GPU Model & NVIDIA GeForce RTX 5070 \\
        Streaming Multiprocessors (SMs) & 48 \\
        Global Memory & 12,226 MiB \\
        Maximum Threads per Block & 1,024 \\
        CUDA Toolkit Version & 13.5 \\
        \bottomrule
    \end{tabular}
\end{table}



\begin{table}[H]
    \centering
    \caption{Dense k-mer frequency array memory requirements
    and feasibility on the NVIDIA GeForce RTX 5070}
    \label{tab:gpu_kmer_memory}
    \begin{tabular}{ccc}
        \toprule
        \textbf{k-mer Length ($k$)} &
        \textbf{Array Memory} &
        \textbf{GPU Feasibility} \\
        \midrule
        12 & 128 MiB & Yes \\
        13 & 512 MiB & Yes \\
        14 & 2 GiB & Yes \\
        15 & 8 GiB & Potentially \\
        16 & 32 GiB & No \\
        \bottomrule
    \end{tabular}
\end{table}

\vspace{\baselineskip}

\FloatBarrier

\section{Execution}

\vspace{\baselineskip}

\subsection{Building and Running}

\vspace{\baselineskip}

\begin{center}
\begin{tcolorbox}[
    colback=gray!8,
    colframe=gray!25,
    arc=3mm,
    boxrule=0.4pt,
    width=0.85\textwidth,
    left=10pt,
    right=10pt,
    top=8pt,
    bottom=8pt
]
\ttfamily
\# Build all implementations

\vspace{6pt} 

make all

\vspace{12pt}

./scripts/benchmark.sh

\vspace{6pt}

\# Sequential
\vspace{6pt}

./sequential/out <k-val> <FASTA\_file> [output\_counts.csv]

\vspace{12pt}

\# OpenMP
\vspace{6pt}

./openMP/out <k> <FASTA\_file> <threads> [output\_counts.csv]

\vspace{12pt}

\# CUDA
\vspace{6pt}

./cuda/out <k> <FASTA\_file> [output\_counts.csv]

\vspace{12pt}

\# Correctness Check
\vspace{6pt}

./scripts/verify\_counts.sh

\end{tcolorbox}
\end{center}

\vspace{\baselineskip}

\section{Experimental Results}

\vspace{\baselineskip}

\subsection{Performance Comparison}

\vspace{\baselineskip}

The sequential, OpenMP, and CUDA implementations were benchmarked using both the E. coli reference
genome and human chromosome 18. Experiments were performed for k-mer lengths of 3, 6, 9, and 12,
with ten trials conducted for each configuration. OpenMP was evaluated using 1, 2, 4, 8, and 16
threads, while the CUDA implementation used a fixed configuration of 256 threads per block.

\vspace{\baselineskip}

Table~\ref{tab:benchmark_comparison} compares the mean counting times for the sequential implementation, the 
fastest measured OpenMP configuration, and CUDA under cold-start and warm-context conditions. The OpenMP 
results represent the thread configuration with the lowest mean counting time among the five configurations 
evaluated (1, 2, 4, 8, and 16 threads). Consequently, the thread count varies between workloads. 

\vspace{\baselineskip}

\FloatBarrier

\input{../source/results/midterm/benchmark_comparison.tex}

\FloatBarrier

\vspace{\baselineskip}

Although increasing the number of threads generally improves parallel counting performance, the additional 
memory allocation and merging overhead can outweigh these benefits, particularly at larger k-mer lengths. At 
$k=12$, each thread requires a private 128 MiB frequency array, resulting in significantly greater memory consumption 
as the thread count increases. This helps explain why configurations with fewer threads achieved better overall counting 
times for some workloads.

\vspace{\baselineskip}

\subsection{CUDA Timing Analysis}

\vspace{\baselineskip}

The CUDA implementation demonstrated a significant difference between cold-start and warm-context performance.
As shown in Table~\ref{tab:benchmark_comparison}, CUDA initialization introduces substantial overhead,
particularly for smaller datasets and k-mer lengths. For the E. coli dataset, the initialization cost
generally exceeded the time required by the sequential implementation to complete the counting operation.

\vspace{\baselineskip}

However, when CUDA initialization is excluded, the GPU implementation becomes more competitive.
For human chromosome 18 at $k=12$, the sequential implementation required an average of 680.26 ms,
while CUDA required 241.53 ms under cold-start conditions and 43.99 ms with an initialized context.
These measurements correspond to speedups of approximately 2.82 and 15.46, respectively.

\vspace{\baselineskip}

To better understand this difference, Table~\ref{tab:cuda_timing} separates CUDA initialization,
kernel execution, warm-context counting, and cold-start counting times.

\vspace{\baselineskip}

\input{../source/results/midterm/cuda_timing_table.tex}

\vspace{\baselineskip}

\subsection{Sequential, OpenMP, and CUDA Performance Comparison}

\vspace{\baselineskip}

The experimental results demonstrate that the effectiveness of parallelization depends on both the
dataset size and k-mer length. For the E. coli dataset, the sequential implementation generally
outperformed OpenMP and cold-start CUDA. The smaller input size limits the benefits of parallel
execution, as the overhead associated with thread management, memory allocation, and GPU
initialization becomes significant relative to the counting workload.

\vspace{\baselineskip}

For human chromosome 18, OpenMP provided improved performance over the sequential implementation
across all tested k-mer lengths. At $k=12$, the sequential implementation required an average of
680.26 ms, while the fastest OpenMP configuration required 381.61 ms using four threads. This
corresponds to a speedup of approximately 1.78. However, the additional memory required for
thread-local counter arrays reduced the benefits of increasing the thread count beyond four.

\vspace{\baselineskip}

CUDA demonstrated the greatest performance improvement for chromosome 18 at $k=12$. The
warm-context implementation achieved a speedup of approximately 15.46 compared with the sequential
baseline, while cold-start CUDA achieved a speedup of approximately 2.82. These results demonstrate
that GPU acceleration can provide substantial improvements for larger counting workloads, but
initialization overhead can significantly affect performance when processing smaller datasets.

\vspace{\baselineskip}

\begin{figure}[H]
    \centering
    \includegraphics[width=0.85\textwidth]
    {../source/results/midterm/figures/chr18/chr18.fa\_speedup.png}
    \caption{OpenMP counting speedup relative to the sequential
    implementation for chromosome 18 across different thread counts
    and k-mer lengths.}
    \label{fig:openmp_speedup}
\end{figure}

\vspace{\baselineskip}

\begin{figure}[H]
    \centering
    \includegraphics[width=0.85\textwidth]
    {../source/results/midterm/figures/chr18/chr18.fa\_cuda\_speedup.png}
    \caption{CUDA counting speedup relative to the sequential
    implementation for chromosome 18, comparing cold-start
    execution with an initialized CUDA context.}
    \label{fig:cuda_speedup}
\end{figure}

\vspace{\baselineskip}

Although the experimental evaluation was limited to k-mer lengths of 12 or less, the CUDA implementation may provide 
additional performance benefits at larger k-mer lengths. As shown in Table~\ref{tab:gpu_kmer_memory}, the NVIDIA 
GeForce RTX 5070 has sufficient global memory to accommodate dense frequency arrays for k-mer lengths of 13 and 14, 
requiring approximately 512 MiB and 2 GiB, respectively. A length of 15 would require approximately 8 GiB, which may 
also be feasible depending on additional GPU memory requirements.

\vspace{\baselineskip}

Increasing k further increases the size of the frequency array, potentially making memory initialization and random-access 
frequency updates more expensive for the sequential implementation. GPU parallelism may help reduce the associated counting 
overhead, although larger arrays also introduce additional memory-bandwidth and allocation costs. Consequently, the performance 
benefits observed at $k=12$ cannot necessarily be extrapolated to larger values without further testing. Evaluating k-mer lengths 
beyond 12 is therefore a potential direction for future testing.

\vspace{\baselineskip}

\section{Conclusion}

This project evaluated the effectiveness of CPU and GPU parallelization for k-mer counting in genomic
sequences. Three implementations were developed using sequential C++, OpenMP, and CUDA, and their
performance was compared across two genomic datasets and four k-mer lengths. All three implementations
were validated against an independent Python reference, confirming the correctness of the frequency
counts for the tested configurations.

\vspace{\baselineskip}

The experimental results demonstrate that parallelization does not always guarantee improved performance.
For smaller datasets, the sequential implementation generally performed better due to the additional
overhead introduced by parallel execution. OpenMP demonstrated improvements on the larger chromosome 18
dataset, although its performance was limited by the memory requirements of thread-local frequency arrays
and the overhead associated with merging their results.

\vspace{\baselineskip}

CUDA demonstrated the greatest performance improvement for larger workloads. For chromosome 18 at
$k=12$, CUDA achieved a speedup of approximately 2.82 under cold-start conditions and 15.46 with an
initialized CUDA context. These results highlight the importance of distinguishing GPU kernel execution
from the additional costs of initialization, memory allocation, and data transfers when evaluating
parallel performance.

\vspace{\baselineskip}

Overall, the results suggest that the effectiveness of parallel k-mer counting depends on the dataset size,
k-mer length, memory requirements, and parallelization strategy. Future work will investigate larger
k-mer lengths, alternative OpenMP memory management approaches, and additional CUDA configurations
to determine whether the observed performance improvements can be maintained as computational and
memory requirements increase.

\vspace{\baselineskip}

%========================================================== Document End

\clearpage

\bibliographystyle{IEEEtran}
\bibliography{references}

\end{document}
%====================================================================================================================== EOF  


